% Part: first-order-logic
% Chapter: proof-systems
% Section: seqeunt-calculus

\documentclass[../../../include/open-logic-section]{subfiles}

\begin{document}

\iftag{FOL}
      {\olfileid{fol}{prf}{seq}}
      {\olfileid{pl}{prf}{seq}}

\olsection{De sequentencalculus}

Veel stelsels van afleidingsregels werken met !!{sentence}s die op een
bepaalde manier bij elkaar staan. De sequentencalculus werkt met
\emph{sequenten}. Een sequent is een uitdrukking van deze vorm:
\[
!A_1, \dots, !A_m \Sequent !B_1, \dots, !B_m,
\]
Links en rechts staat een rij !!{sentence}s. Het sequententeken
$\Sequent$ scheidt de twee rijen van elkaar. Elk van beide rijen mag
leeg zijn. !!^a{derivation} in de sequentencalculus vormt een boom van
sequenten. De sequenten bovenaan hebben een bijzondere vorm. Ze heten
``beginsequenten'' of ``axioma's''. Elke andere sequent volgt met een
afleidingsregel uit de sequenten die er direct boven staan. Zo'n regel
kan links of rechts !!{sentence}s toevoegen, weghalen of verplaatsen.
Een regel kan ook in zijn conclusie een samengestelde !!{formula}
invoeren. Met $\LeftR{\land}$ mogen we bijvoorbeeld van $!A, \Gamma
\Sequent \Delta$ naar $!A \land !B, \Gamma \Sequent \Delta$ gaan. Met
$\RightR{\lif}$ mogen we van $!A, \Gamma \Sequent \Delta, !B$ naar
$\Gamma \Sequent \Delta, !A \lif !B$ gaan. Dit geldt voor alle
$\Gamma$, $\Delta$, $!A$ en~$!B$. (Daarbij mogen ook $\Gamma$
en~$\Delta$ leeg zijn.)

Voor de sequentencalculus definiëren we de relatie $\Proves$ als
volgt. $\Gamma \Proves !A$ geldt precies als er een rij $\Gamma_0$
bestaat met twee eigenschappen. Elke $!A$ in $\Gamma_0$ behoort
tot~$\Gamma$, en er bestaat een !!{derivation} met de
sequent~$\Gamma_0 \Sequent !A$ onderaan, aan de wortel van de boom.
$!A$ is een stelling in de sequentencalculus als de
sequent~$\Sequent !A$ !!a{derivation} heeft. De volgende
!!a{derivation} laat bijvoorbeeld zien dat $\Proves (!A \land !B) \lif
!A$:
\begin{prooftree}
  \Axiom$!A \fCenter !A$
  \RightLabel{\LeftR{\land}}
  \UnaryInf$!A \land !B \fCenter !A$
  \RightLabel{\RightR{\lif}}
  \UnaryInf$\fCenter (!A \land !B) \lif !A$
\end{prooftree}

Een verzameling $\Gamma$ is inconsistent in de sequentencalculus als
er !!a{derivation} van $\Gamma_0 \Sequent$ bestaat. Daarbij behoort
elke $!A \in \Gamma_0$ tot~$\Gamma$ en is de rechterkant van de
sequent leeg. Met de regel \RightR{\Weakening} kan uit een
inconsistente verzameling elke !!{sentence} worden !!{derive}d.

Gerhard Gentzen vond de sequentencalculus in de jaren dertig van de
twintigste eeuw uit. De opzet is systematisch en symmetrisch. Daardoor
is dit systeem heel geschikt om een theorie over !!{derivation}s op te
bouwen.
Het is vrij gemakkelijk om in de sequentencalculus een
!!{derivation} te vinden. Zo'n !!{derivation} is vaak wel moeilijk te
lezen. Ook is niet altijd duidelijk hoe hij samenhangt met een bewijs.
Toch is de sequentencalculus een zeer elegante aanpak gebleken, en voor
veel logica's bestaat een sequentencalculus.
\end{document}
